\documentclass[11pt,a4paper]{article}
\usepackage[T1]{fontenc}
\usepackage{lmodern,amsmath,amssymb,amsthm,mathtools,microtype}
\usepackage[margin=27mm]{geometry}
\usepackage[colorlinks=true,linkcolor=blue,citecolor=blue,urlcolor=blue]{hyperref}
\newtheorem{theorem}{Theorem}[section]
\newtheorem{lemma}[theorem]{Lemma}
\newtheorem{proposition}[theorem]{Proposition}
\newtheorem{corollary}[theorem]{Corollary}
\newcommand{\Q}{\mathbb Q}\newcommand{\Z}{\mathbb Z}
\newcommand{\OK}{\mathcal O_K}
\newcommand{\hfin}{h_{\mathrm{fin}}}
\newcommand{\hinf}{h_{\infty}}
\DeclareMathOperator{\ord}{ord}
\setlength{\emergencystretch}{2em}
\title{Effective heights of sums on cyclic elliptic groups over number fields}
\author{}\date{17 November 2026}
\begin{document}\maketitle
\begin{abstract}
Let $E$ be an elliptic curve over a number field and let $P$ have
infinite order. The height of a sum of finitely many abscissas of
points $[n_i]P+T_i$, with rational torsion translates $T_i$, is at
least the largest summand height minus an effective logarithmic
index term. The result holds over arbitrary number fields and for
every elliptic curve. The finite-place estimate has explicit
constants for any integral short Weierstrass model. The
archimedean estimate follows from the optimal coefficient-height
dependence in the theorem of David and Hirata--Kohno.
An effective finiteness statement for integral and $S$-integral
abscissa sums follows. These results concern a fixed cyclic group
modulo torsion; neither field traces nor arbitrary collections of
elliptic points satisfy the same assertion.
\end{abstract}

\section{Statement and normalization}
Let $K$ be a number field of degree $d_K$ and let
\[
 E:\quad y^2=x^3+Ax+B,\qquad A,B\in\OK,\qquad 4A^3+27B^2\ne0.
\]
Fix an infinite-order point $P\in E(K)$ and an integer $r\ge1$.
For $1\le i\le r$ take
\[
 Q_i=[n_i]P+T_i\ne O,\qquad n_i\in\Z,\qquad T_i\in E(K)_{\rm tors}.
\]
Put $M=\max(1,|n_1|,\ldots,|n_r|)$. Repeated points, negative
indices and index zero are allowed, provided every $Q_i$ is finite.
We write $h$ for the absolute logarithmic height and decompose it
as $h=\hfin+\hinf$. At a finite place $\mathfrak p\mid p$,
$v_{\mathfrak p}$ has value group $\Z$; the local contribution is
\[
 \frac{f_{\mathfrak p}\log p}{d_K}
                    \max(0,-v_{\mathfrak p}(a)).
\]
All these heights of zero are zero. Define
\[
 a_r=(r-1)2^{r-1},\qquad
 b_r=\sum_{j=1}^r\binom{r-1}{j-1}\log j,
\]
and
\[
 R_K=\frac{2}{d_K}
 \sum_{\substack{\mathfrak p\mid p\\p\ne2}}
 f_{\mathfrak p}\left\lfloor\frac{e_{\mathfrak p}-1}{2}\right\rfloor\log p.
\]
Only odd ramified places with ramification index at least three
contribute to $R_K$.

\begin{theorem}\label{thm:main}
There is an effectively computable constant $C=C(E,K,P)>0$ such
that, for every tuple as above and every $c\in K$, setting
$Z_c=\sum_i x(Q_i)+c$ gives
\begin{align}
 \hfin(Z_c)&\ge\max_i\hfin(x(Q_i))
       -a_r\log M-b_r-R_K-\hfin(c),\label{eq:finite}\\
 h(Z_c)&\ge\max_i h(x(Q_i))
       -a_r\log M-b_r-R_K-C\log(M+2)-\hfin(c).
       \label{eq:global}
\end{align}
Moreover
\[
 R_K\le\frac{\log|\operatorname{disc}(K)|}{d_K}.
\]
The constant $C$ is independent of $r,c,n_i$ and the chosen
$K$-rational torsion translates. In particular, for fixed $r$
and integral $c$ the loss in~\eqref{eq:global} is $O_{E,K,P,r}(\log(M+2))$,
with an effective implied constant.
\end{theorem}

Every elliptic curve over a number field admits an integral short
Weierstrass model after a change of variables and clearing fixed
denominators. Thus the theorem applies to every elliptic curve;
changing its fixed abscissa changes only effective constants.
The explicit finite-place bound refers to the displayed integral
model. No assertion is made that the David--Hirata--Kohno constant
has a convenient numerical size.

\section{Integral formal expansions}
Write $z=-x/y$ at the identity and $X=z^2x$. The curve equation
becomes
\[
 X^3-X^2+Az^4X+Bz^6=0.
\]
Its derivative with respect to $X$ is one at $(X,z)=(1,0)$.
Coefficient recursion therefore gives
\[
 X\in\Z[A,B][[z^2]],\qquad X=1+O(z^4).
\]
In particular $x=z^{-2}+O(z^2)$, with integral coefficients.
The invariant differential satisfies
\[
 \frac{dx}{2y}=\left(1-\frac{zX'}{2X}\right)dz.
\]
Since $X$ is an even integral series and $X(0)=1$, the factor in
parentheses belongs to $1+z^4\Z[A,B][[z^2]]$. Its integral with
zero constant term is the formal logarithm
\[
 \ell(z)=z+\sum_{\substack{m\ge5\\m\text{ odd}}}\frac{a_m}{m}z^m,
 \qquad a_m\in\OK.
 \label{eq:log}
\]
These calculations do not require good reduction.

At $\mathfrak p\mid p$, put $e=e_{\mathfrak p}$ and call a formal
point \emph{deep} if its depth $v_{\mathfrak p}(z)$ is at least
$\lceil e/2\rceil$ when $p$ is odd, or at least one when $p=2$.
The deep points together with $O$ form an open subgroup of
$E(K_{\mathfrak p})$: the integral formal group law and its inverse
preserve every positive valuation ideal. Set
\[
 b_{\mathfrak p}=
 \begin{cases}2\lfloor(e-1)/2\rfloor,&p\ne2,\\0,&p=2.\end{cases}
\]
A finite point outside the deep subgroup has
$v_{\mathfrak p}(x)\ge-b_{\mathfrak p}$. Indeed, a point of negative
abscissa valuation is formal, and there
$v_{\mathfrak p}(x)=-2v_{\mathfrak p}(z)$.

\begin{lemma}\label{lem:principal}
For every nonidentity deep point,
\[
 v_{\mathfrak p}(\ell(z))=v_{\mathfrak p}(z),\qquad
                 x(z)-\ell(z)^{-2}\in\mathcal O_{K_{\mathfrak p}}.
\]
The formal logarithm is injective on the deep subgroup.
\end{lemma}
\begin{proof}
Put $a=v_{\mathfrak p}(z)$. For every odd $m\ge5$ and odd prime
$p$ we have $v_p(m)\le(m-3)/2$. To see this, if the reverse
integer inequality held, then $m\ge3^{(m-1)/2}>m$; the latter
strict inequality holds for odd $m\ge5$. When $2a\ge e$,
\[
 ma-e v_p(m)\ge ma-2a(m-3)/2=3a.
\]
The valuations tend to infinity, since $v_p(m)=O(\log m)$.
At $p=2$ every denominator $m$ is a unit, and every term beyond
$z$ has valuation at least $5a$. Thus in all cases
$\ell(z)=z(1+\eta)$ with $v_{\mathfrak p}(\eta)\ge2a$.
It follows that $\ell(z)^{-2}-z^{-2}$ is integral, as is
$x(z)-z^{-2}$. The logarithm is a homomorphism on the formal
group, and its valuation identity gives injectivity on this
subgroup.
\end{proof}

Compactness of $E(K_{\mathfrak p})$ implies that a positive multiple
$[q]P$ belongs to the deep subgroup. Put
$L_{\mathfrak p}=\ell([q]P)/q$. The homomorphism law makes this
independent of $q$. It is nonzero because $P$ is non-torsion and
the logarithm is injective on the deep subgroup. If $Q_i$ is deep,
choose a common multiple killing $T_i$ and carrying $P$ into that
subgroup. Applying the logarithm gives
\[
 \ell(Q_i)=n_iL_{\mathfrak p}.
\]
Consequently a deep $Q_i$ has $n_i\ne0$, and
\[
 x(Q_i)=(n_iL_{\mathfrak p})^{-2}+\varepsilon_i,
             \qquad v_{\mathfrak p}(\varepsilon_i)\ge0.\tag{*}
\]
This also handles torsion translates of every order: no assertion
that a shallow torsion abscissa is integral is used.

\section{The finite-place inequality}
Fix $i$. At a place where $Q_i$ is deep, let $J$ be the set of
all deep points among the $Q_j$ and put
\[
 C_J=\sum_{j\in J}\prod_{a\in J\setminus\{j\}}n_a^2\in\Z_{>0}.
\]
All the indices in $J$ are nonzero. Positivity ensures that $C_J$
cannot vanish, even for negative indices. Summing~$(*)$ gives
\[
 Z_c=(n_iL_{\mathfrak p})^{-2}
       \frac{C_J}{\prod_{j\in J\setminus\{i\}}n_j^2}+\varepsilon,
\]
where
\[
 v_{\mathfrak p}(\varepsilon)\ge
 -b_{\mathfrak p}-\max(0,-v_{\mathfrak p}(c)).
\]
Write $H=-v_{\mathfrak p}(x(Q_i))$ and let $A_J$ denote the
rational quotient in the principal term. If
$H-v_{\mathfrak p}(A_J)$ exceeds the displayed error allowance,
the principal term controls the valuation of $Z_c$. Otherwise
the entire local height of $x(Q_i)$ is at most
$v_{\mathfrak p}(A_J)$ plus that allowance. In both cases the
local height loss, before multiplying by its place weight, is at
most
\[
 v_{\mathfrak p}(C_J)+b_{\mathfrak p}
                        +\max(0,-v_{\mathfrak p}(c)),
\]
because the denominator of $A_J$ is an integer. If $Q_i$ is
not deep, its whole local height is at most $b_{\mathfrak p}$.

For each fixed positive integer $C$, the sum over all finite
places of its weighted valuation is $\log C$. Each possible
subset $J$ containing $i$ satisfies
$C_J\le |J|M^{2(|J|-1)}$. Summing over the at most
$\binom{r-1}{j-1}$ subsets of size $j$ gives
\[
 \sum_{j=1}^r\binom{r-1}{j-1}
       \bigl(\log j+2(j-1)\log M\bigr)
                  =b_r+a_r\log M.
\]
The remaining local allowances sum to $R_K+\hfin(c)$. This
proves~\eqref{eq:finite} for every $i$, hence for their maximum.
The argument also covers $n_i=0$, since that point is never deep.

Finally $2\lfloor(e-1)/2\rfloor\le e-1$, whereas the exponent
of the local different is at least $e-1$. The discriminant formula
therefore gives the asserted bound on $R_K$.

\section{The archimedean estimate}
The coefficient-height dependence in the following result is
essential. Theorem 1.6 of David and Hirata--Kohno~\cite{DHK}
implies the fixed-data statement used here:

\begin{lemma}\label{lem:dhk}
Fix elliptic curves over a number field and fixed elliptic
logarithms $u_1,\ldots,u_k$ of algebraic points, where periods
are permitted as logarithms of the identity. There is an
effectively computable $C_0>0$ such that for any integers
$b_0,\ldots,b_k$, if
$\Lambda=b_0+\sum_{j=1}^k b_j u_j\ne0$, then
\[
 \log|\Lambda|\ge-C_0\log\bigl(2+\max_j|b_j|\bigr).
\]
\end{lemma}
\begin{proof}[Verification of the specialization]
Use Theorem 1.6 with one linear form, and use the fixed vector
$u=(1,u_1,\ldots,u_k)$. Since $\Lambda\ne0$, this vector is not
in the kernel of the form, and cannot belong to the tangent
space of a subgroup contained in that kernel. This is precisely
the theorem's nonvanishing hypothesis $s=t=1$.
In Data 1.4, choose the auxiliary parameter $\mathcal E=e$ and
choose each $V_j$ once and for all large enough to satisfy the
canonical-height and analytic-size inequalities. This is possible
because the points, logarithms, periods and number field are
fixed. The degree, model heights, modular heights and all $V_j$
are then fixed. Only the coefficient-height parameter $B_0$
varies, and one can take $B_0=\max(e,|b_0|,\ldots,|b_k|)$.
In the stated bound, $\log B_0$ occurs linearly; the iterated
logarithm is $\log\log^+ V$, with $V=\max V_j$, not an iterated
logarithm of $B_0$. All remaining factors are fixed and effective.
This yields the displayed inequality after enlarging $C_0$.
The paper's projective exponential map includes the identity,
so no exclusion of period logarithms is needed.
\end{proof}

\begin{proposition}\label{prop:arch}
There is an effectively computable $C=C(E,K,P)>0$ such that
\[
 \hinf\bigl(x([n]P+T)\bigr)\le C\log(2+|n|)
\]
for every $n\in\Z$, every $T\in E(K)_{\rm tors}$, and every
finite $[n]P+T$.
\end{proposition}
\begin{proof}
Fix an embedding $\sigma:K\hookrightarrow\mathbb C$. For the
complex uniformization of the displayed curve, choose periods
$\omega_1,\omega_2$ and an elliptic logarithm $u_P$ in their
closed centered fundamental parallelogram. The coordinate map is
$x=\wp(u)$ and $y=\wp'(u)/2$, with invariants
$g_2=-4\sigma(A)$ and $g_3=-4\sigma(B)$.
Choose an effectively known positive integer $t$ annihilating
$E(K)_{\rm tors}$. Such an integer can, for example, be obtained
from reduction counts at two good places of distinct residue
characteristics. If their counts are $a,b$, then $t=ab$ suffices:
for every rational prime $\ell$ at least one reduction is injective
on the $\ell$-primary torsion, and its order divides the corresponding
count.

For a torsion point $T$, choose
$u_T=(a_T\omega_1+b_T\omega_2)/t$, with
$0\le a_T,b_T<t$. Reduce $nu_P+u_T$ to a representative $u_Q$
in the centered parallelogram. Then
\[
 t u_Q=tnu_P+(a_T-tm_1)\omega_1+(b_T-tm_2)\omega_2
\]
for integers $m_1,m_2$ satisfying
$|m_j|\le |n|/2+2$. Thus every coefficient on the right has
absolute value at most $t(|n|+3)$. The right side is nonzero,
since $Q=[n]P+T\ne O$ and the only lattice point in this
parallelogram is zero.

Apply Lemma~\ref{lem:dhk} to the three fixed logarithms
$u_P,\omega_1,\omega_2$ of $P,O,O$. It gives
$-\log|u_Q|\le C_1\log(2+|n|)$ with an effective constant.
On the fixed closed parallelogram, $|u|^2|\wp(u)|$ is bounded,
with its removable value at zero defined by the Laurent
expansion. Its bound is effective from the convergent
Weierstrass series, or from the analytic estimates in
\cite[Section 2.1]{DHK}. Hence
\[
 \log\max(1,|\sigma(x(Q))|)\le C_2\log(2+|n|).
\]
This proof also covers $n=0$: the linear form is still nonzero
for a finite torsion point. Summing with the normalized
archimedean weights proves the proposition.
\end{proof}

Since $h(Z_c)\ge\hfin(Z_c)$, applying Proposition~\ref{prop:arch}
to an index maximizing $h(x(Q_i))$ and then using~\eqref{eq:finite}
proves~\eqref{eq:global}. This completes the proof of
Theorem~\ref{thm:main}.

\section{Effective finiteness and the role of the cyclic group}
Normalize the canonical height by
\[
 \widehat h(Q)=\tfrac12\lim_{k\to\infty}4^{-k}h(x([2^k]Q)).
\]
The usual effective comparison with naive height gives
$h(x(Q))=2\widehat h(Q)+O_E(1)$; see~\cite[Chapter VIII,
Section 9]{Silverman}. The constant can be obtained from
resultant bounds for the duplication morphism.

\begin{corollary}\label{cor:growth}
For fixed $r$ and integral $c$, there is an effective constant
$C_r$ such that
\[
 h\left(\sum_i x(Q_i)+c\right)
 \ge 2\widehat h(P)\max_i n_i^2-C_r\log(M+2).
\]
More strongly, for any fixed finite set $S$ of finite places,
there is an effective $C_{r,S}$ such that the same lower bound,
with $C_r$ replaced by $C_{r,S}$, holds for the height contribution
outside $S$ whenever $c$ is $S$-integral. Consequently there are
only finitely many tuples of indices and rational torsion
translates for which $\sum_i x(Q_i)$ is $S$-integral, and all
indices are effectively bounded.
\end{corollary}
\begin{proof}
The first assertion follows from Theorem~\ref{thm:main} and
$\widehat h([n]P+T)=n^2\widehat h(P)$.
For the stronger assertion, the finite-place proof can simply be
summed over places outside $S$. At a fixed place in $S$, a deep
point satisfies
\[
 \max(0,-v_{\mathfrak p}(x([n]P+T)))
       =2v_{\mathfrak p}(n)+2v_{\mathfrak p}(L_{\mathfrak p}),
\]
where $n\ne0$. The first term is at most
$2e_{\mathfrak p}\log M/\log p$; the second is fixed and
computable from any sufficiently deep multiple of $P$. A
nondeep point has bounded local height. Thus the sum of local
heights over $S$, like the archimedean contribution, is
$O_{E,K,P,S}(\log(M+2))$, effectively. Subtracting these
contributions from the naive height and applying the
finite-place inequality outside $S$ proves the claim.
If the sum is $S$-integral, this outside height is zero, whereas
$\widehat h(P)>0$. A positive quadratic eventually exceeds every
fixed logarithmic term, which gives an effective index bound.
There are only finitely many bounded integer tuples and
$K$-rational torsion points.
\end{proof}

\begin{corollary}\label{cor:rankonecurves}
For every number field $K$ there exist infinitely many elliptic
curves $E/K$, up to isomorphism, for which the following holds.
Choose an integral short Weierstrass model and fix $r\ge1$.
There is a constant $C_r>0$ such that, for all finite points
$Q_1,\ldots,Q_r\in E(K)$ and all $c\in\OK$, putting
$H=\max_i h(x(Q_i))$ gives
\[
 h\left(\sum_i x(Q_i)+c\right)\ge H-C_r\log(H+2).
\]
For a fixed curve and a supplied generator modulo torsion,
$C_r$ is effectively computable. For each fixed finite set $S$,
only finitely many tuples of finite $K$-rational points have an
$S$-integral abscissa sum.
\end{corollary}
\begin{proof}
Zywina's rank-one existence theorem~\cite[Theorem 1.1]{Zywina}
supplies infinitely many such curves. On each one, choose $P$
generating the free part of $E(K)$. Every finite point is then
$[n]P+T$ with $T$ rational torsion. Canonical-height comparison
gives $\max_i n_i^2\le C_E(H+1)$, so
$\log(M+2)\le C'_E\log(H+2)$. Apply
Theorem~\ref{thm:main} and Corollary~\ref{cor:growth}.
\end{proof}

The use of a common cyclic group modulo torsion is indispensable
in the finite-place proof: it produces the positive integer
$C_J$. For the curve $y^2=x^3-Dx$ over a field containing $i$,
the automorphism $\iota(x,y)=(-x,iy)$ gives
$x(Q)+x(\iota Q)=0$. If $Q$ varies through non-torsion multiples,
its abscissa height is unbounded. The two families generally
span two independent directions over $\Z$, so they are outside
the theorem's hypothesis. Likewise a field trace can cancel
between conjugates that do not lie in the same cyclic group
modulo torsion. These obstructions are developed in~\cite{CM}.

The theorems above do not supply a Diophantine definition of a
ring of integers in its fraction field. Such an application
would still need a Diophantine description of the required
cyclic point domain, an integrality detector, and complete
rational witness constructions. Replacing a sum by a field trace
is an additional step not justified by the present estimates.

\begin{thebibliography}{9}
\bibitem{DHK}
S. David and N. Hirata--Kohno,
\emph{Linear forms in elliptic logarithms},
Journal f\"ur die reine und angewandte Mathematik \textbf{628}
(2009), 37--89. Data 1.4 and Theorems 1.5--1.6, pp.~41--43.
\url{https://doi.org/10.1515/CRELLE.2009.018}.
\bibitem{Silverman}
J. H. Silverman, \emph{The Arithmetic of Elliptic Curves},
second edition, Graduate Texts in Mathematics \textbf{106},
Springer, 2009.
\bibitem{Zywina}
D. Zywina, \emph{Rank one elliptic curves and rank stability},
preprint, 2025, arXiv:2505.16960v1, Theorem 1.1.
\url{https://arxiv.org/abs/2505.16960}.
\bibitem{CM}
\begingroup\raggedright
\emph{Complex Multiplication and Cancellation in Elliptic Trace Heights},
10 November 2026,
\nolinkurl{negative_results/papers/number_fields/complex_multiplication_and_trace_cancellation_20261110_v19.tex}.\par\endgroup
\end{thebibliography}
\end{document}
